¿Cuál es la diferencia principal? A menudo se usan como sinónimos, pero en la práctica son dos conceptos distintos que trabajan juntos:

\begin{itemize}
    \item \textbf{El Algoritmo (El Motor / La Receta):} Es el procedimiento matemático o el conjunto de reglas que la computadora utiliza para analizar los datos, encontrar patrones y ``aprender'' de ellos. Ejemplos de algoritmos son el Descenso del Gradiente o la Retropropagación (\textit{Backpropagation}).
    \item \textbf{El Modelo (El Resultado / El Plato Terminado):} Es la representación matemática final que resulta \textbf{después} de que el algoritmo ha procesado los datos de entrenamiento. Es esencialmente una función matemática a la que le entregas nuevos datos de entrada para obtener una predicción o clasificación de salida.
\end{itemize}



En resumen: \textbf{Algoritmo + Datos de Entrenamiento = Modelo de Machine Learning.}

\begin{figure}[htbp]
    \centering
    \shorthandoff{>}
    % Ajusta todo el contenido al ancho exacto de la línea de texto (\textwidth)
    \resizebox{\textwidth}{!}{%
    \begin{tikzpicture}[
        % Definición de estilos para las cajas (tamaños ligeramente reducidos)
        base/.style = {rectangle, rounded corners, draw=black, minimum width=3cm, minimum height=1.5cm, text centered, font=\sffamily, align=center},
        data/.style = {base, fill=blue!10, draw=blue!80!black, thick},
        algo/.style = {base, fill=orange!10, draw=orange!80!black, thick},
        model/.style = {base, fill=green!10, draw=green!80!black, thick},
        param/.style = {base, fill=yellow!15, draw=yellow!80!black, dashed, minimum height=1cm},
        arrow/.style = {thick, ->, >=stealth}
        ]

        % 1. Nodos Principales (distancia reducida a 0.8cm)
        \node (data) [data] {\textbf{1. Datos} \\ (Training Data) \\ $X, y$};
        
        \node (algo) [algo, right=0.8cm of data] {\textbf{2. Algoritmo} \\ Optimización \\ (Ej: Gradiente)};
        
        \node (model) [model, right=0.8cm of algo] {\textbf{3. Modelo} \\ Función \\ $f(x)$};
        
        \node (pred) [base, fill=gray!10, right=0.8cm of model] {\textbf{4. Predicción} \\ Salida \\ $\hat{y}$};

        % 2. Nodos Secundarios
        \node (hyper) [param, below=0.8cm of algo] {\textbf{Hiperparámetros} \\ (Configuración)};
        
        \node (params) [param, below=0.8cm of model] {\textbf{Parámetros} \\ Pesos ($w$), Sesgos ($b$)};

        % 3. Conexiones (Flechas)
        \draw [arrow] (data) -- node[above, font=\footnotesize] {Alimenta} (algo);
        \draw [arrow] (algo) -- node[above, font=\footnotesize] {Genera} (model);
        \draw [arrow] (model) -- node[above, font=\footnotesize] {Datos} (pred);
        
        \draw [thick, dashed, ->, >=stealth] (hyper) -- (algo);
        \draw [thick, dashed, <-, >=stealth] (params) -- (model);

        % 4. Caja de fondo para agrupar la fase de aprendizaje
        \begin{scope}[on background layer]
            \node[draw=gray!50, dashed, rounded corners, fit=(algo)(model)(params)(hyper), inner sep=0.4cm, fill=gray!5] (box) {};
            \node[above=0.2cm of box.north, font=\bfseries\sffamily, text=gray!70!black] {Fase de Aprendizaje};
        \end{scope}

    \end{tikzpicture}%
    } % <-- Fin del resizebox
    \caption{Visualización del flujo: Cómo los algoritmos utilizan los datos para ajustar parámetros y generar un modelo predictivo.}
    \label{fig:ml_flow}
\end{figure}

%\vspace{1.4cm} % Añade un pequeño espacio vertical de separación con el gráfico
\noindent \textbf{Explicación del flujo de Machine Learning:}
Este diagrama ilustra la transición desde la información sin procesar hasta la toma de decisiones automatizada:

\begin{description}
    \item[1. El punto de partida:] El proceso arranca cuando los \textbf{Datos de Entrada} alimentan al \textbf{Algoritmo}. El algoritmo actúa como un motor de optimización que busca patrones matemáticos. Para que este motor funcione adecuadamente, el científico de datos establece reglas previas llamadas \textit{Hiperparámetros} (configuraciones manuales que guían el aprendizaje).
    
    \item[2. La consolidación (El Aprendizaje):] Al procesar los datos, el algoritmo ajusta sus valores internos. El resultado final de esta fase es el \textbf{Modelo}. Este modelo almacena los \textit{Parámetros Aprendidos} (como los pesos y sesgos), que son los valores exactos que el algoritmo descubrió por sí solo para resolver el problema.
    
    \item[3. La fase de inferencia:] Una vez que la etapa de aprendizaje concluye, el algoritmo pasa a un segundo plano. Ahora es el \textbf{Modelo}, con sus parámetros ya calibrados, el que recibe datos nuevos que nunca antes ha visto para aplicar su función matemática y generar una \textbf{Predicción} final.
\end{description}

\subsubsection{Ejemplos de Algoritmos y sus Modelos}

A continuación, se presentan algunos de los enfoques más clásicos en Machine Learning, cómo funciona su algoritmo y la estructura matemática del modelo resultante.

\begin{enumerate}
    \item {\bf Regresión Lineal (Linear Regression):} 
    Se utiliza para predecir un valor numérico continuo (por ejemplo, el precio de una casa basado en su tamaño).

    \begin{itemize}
        \item \textbf{El Algoritmo:} Usualmente se emplea el método de \textit{Mínimos Cuadrados Ordinarios (OLS)} o el \textit{Descenso del Gradiente} para encontrar los coeficientes que minimicen el error cuadrático entre la predicción y los datos reales.
        \item \textbf{El Modelo:} Una vez entrenado, el modelo es una ecuación lineal (una línea recta o un hiperplano) donde se han calculado los pesos ideales ($\beta$).
    \end{itemize}
    \begin{equation}
        y = \beta_0 + \beta_1 x_1 + \beta_2 x_2 + \dots + \beta_n x_n
    \end{equation}
    \textit{(Donde $y$ es la predicción, $x_i$ son las características de entrada, $\beta_0$ es el sesgo y $\beta_i$ son los pesos aprendidos).}
    \item {\bf Regresión Logística (Logistic Regression)}
A pesar de su nombre, se utiliza para \textbf{clasificación} binaria (por ejemplo, decidir si un correo es ``Spam'' o ``No Spam'').

\begin{itemize}
    \item \textbf{El Algoritmo:} Utiliza la \textit{Maximización de la Verosimilitud} (Maximum Likelihood Estimation) para ajustar los pesos del modelo a los datos de entrenamiento.
    \item \textbf{El Modelo:} Toma una ecuación lineal y la pasa por una función Sigmoide para comprimir el resultado y convertirlo en una probabilidad matemática entre 0 y 1.
\end{itemize}
\begin{equation}
    P(y=1|x) = \frac{1}{1 + e^{-(\beta_0 + \beta_1 x)}}
\end{equation}
\textit{(Donde $P(y=1|x)$ es la probabilidad de que la entrada $x$ pertenezca a la clase 1).}

\item {\bf Máquinas de Vectores de Soporte (SVM)}
Se utiliza principalmente para clasificación compleja. Busca trazar la mejor ``frontera'' posible entre diferentes tipos de datos.

\begin{itemize}
    \item \textbf{El Algoritmo:} Es un solucionador de optimización matemática que busca maximizar el ``margen'' (la distancia) entre los puntos de datos más cercanos de diferentes clases (los vectores de soporte). El objetivo del algoritmo es minimizar $\frac{1}{2} ||w||^2$.
    \item \textbf{El Modelo:} El resultado es la ecuación del \textbf{hiperplano óptimo} que separa las clases en el espacio multidimensional.
\end{itemize}
\begin{equation}
    w^T x + b = 0
\end{equation}
\textit{(Donde $w$ es el vector de pesos que determina la orientación del hiperplano, $x$ es la entrada y $b$ es el sesgo).}

\item {\bf Redes Neuronales Artificiales (ANN)}
Es la base del \textit{Deep Learning}. Modela relaciones altamente complejas y no lineales.

\begin{itemize}
    \item \textbf{El Algoritmo:} Utiliza \textit{Backpropagation} (Propagación hacia atrás) junto con un optimizador para ajustar millones de pesos en la red calculando derivadas parciales (gradientes) de la función de pérdida.
    \item \textbf{El Modelo:} Es una red de funciones matemáticas interconectadas. Para una sola ``neurona'' dentro de la red, el modelo matemático calcula una suma ponderada y la pasa por una función de activación no lineal $\sigma$ (como ReLU o Sigmoide).
\end{itemize}
\begin{equation}
    a = \sigma\left(\sum_{i=1}^{n} w_i x_i + b\right)
\end{equation}
\textit{(Donde $a$ es la salida o activación de la neurona, $w_i$ son los pesos, $x_i$ son las entradas y $b$ es el sesgo).}

\item {\bf Clasificador Naïve Bayes (Naïve Bayes)}
Este es un modelo generativo fundamentado en el Teorema de Bayes, utilizado ampliamente en clasificación de textos y análisis de sentimientos. Se le llama ``ingenuo'' (\textit{naïve}) porque asume independencia condicional estricta entre las características dado la clase.

\begin{itemize}
    \item \textbf{El Algoritmo:} Utiliza la estimación de Máxima Probabilidad a Posteriori (MAP) contando frecuencias en los datos de entrenamiento para calcular empíricamente la probabilidad a priori de cada clase $P(y)$ y las probabilidades condicionales de cada característica $P(x_i | y)$.
    \item \textbf{El Modelo:} Es una regla de decisión matemática. Para una nueva observación con características $x = (x_1, \dots, x_n)$, el modelo asigna la clase $\hat{y}$ que maximiza la probabilidad conjunta (gracias al supuesto de independencia):
\end{itemize}
\begin{equation}
    \hat{y} = \arg \max_{y} P(y) \prod_{i=1}^{n} P(x_i \mid y)
\end{equation}

\item  {\bf Agrupamiento K-Medias (K-Means Clustering)}
Es el algoritmo por excelencia para el aprendizaje no supervisado. Su objetivo es particionar un conjunto de $n$ observaciones en $k$ grupos (clústeres), donde cada observación pertenece al grupo cuyo valor medio sea más cercano.

\begin{itemize}
    \item \textbf{El Algoritmo:} Típicamente emplea el \textit{Algoritmo de Lloyd} (un enfoque heurístico iterativo de Expectativa-Maximización). En cada iteración, reasigna los puntos al centroide más cercano y luego recalcula la posición del centroide como el centro de masa geométrico del grupo. El objetivo del algoritmo es minimizar la varianza intra-clúster (la suma de los errores cuadráticos):
    \begin{equation}
        J = \sum_{j=1}^{k} \sum_{x_i \in C_j} || x_i - \mu_j ||^2
    \end{equation}
    \item \textbf{El Modelo:} Una vez que converge, el modelo no es una ecuación continua, sino un conjunto discreto de $k$ vectores en $\mathbb{R}^d$ que representan los centroides finales $\{ \mu_1, \mu_2, \dots, \mu_k \}$. Para predecir a qué grupo pertenece un nuevo punto $x$, el modelo evalúa la métrica de distancia:
\end{itemize}
\begin{equation}
    \text{Clúster}(x) = \arg \min_{j \in \{1, \dots, k\}} || x - \mu_j ||^2
\end{equation}

\item {\bf Análisis de Componentes Principales (PCA)}
Aunque a menudo se clasifica como una técnica de preprocesamiento, es un modelo matemático de aprendizaje no supervisado utilizado para la reducción de dimensionalidad, preservando la mayor cantidad de varianza posible.

\begin{itemize}
    \item \textbf{El Algoritmo:} Implica álgebra lineal pura. Toma la matriz de datos centrados $X$, calcula su matriz de covarianza $\Sigma = \frac{1}{n-1} X^T X$ y luego realiza una descomposición espectral para encontrar sus autovalores ($\lambda$) y autovectores ($v$). Se seleccionan los $k$ autovectores correspondientes a los $k$ autovalores más grandes.
    \begin{equation}
        \Sigma v = \lambda v
    \end{equation}
    \item \textbf{El Modelo:} Es una matriz de proyección ortogonal $W_{d \times k}$ (cuyas columnas son los $k$ autovectores principales). Al recibir un nuevo dato original $x$ (de dimensión $d$), el modelo lo transforma a un nuevo espacio latente $z$ (de dimensión $k$):
\end{itemize}
\begin{equation}
    z = x W
\end{equation}

\end{enumerate}